\chapter[Systems, gauge, variable coefficients]{Phases 6--8: systems, gauge and variable coefficients (other results)}\label{ch:cIV1}
\chapterbrief
{the notation of Chapters~\ref{ch:cIII1} and~\ref{ch:cIII2}.}
{to give the remaining results of the system and variable-coefficient packages: short proofs in full, and the others as labelled statements with a pointer to the research file.}
{Dim.~B73--B80, and the list cards of this chapter. They prove the corresponding parts of Chapters~\ref{PAP-ch:systems} and~\ref{PAP-ch:varcoef} of the paper.}

\section{Systems}

\begin{dimcard}{B73}{SYS-R1}{Finite data and ODE systems}{\proved}{thm:sysr1}
\textbf{(SYS-R1)} If $u\in\EP^k$ then $M_u$ has finite length. Hence from exponential-polynomial data one identifies only systems with finite-dimensional solution space.

\textbf{(SYS-T0)} For $n=1$, $R=\C[s]$, $P\in R^{l\times k}$ of rank $r$ with invariant factors $e_1\mid\dots\mid e_r$ on an open interval: an identifying $u$ exists iff $e_1=\dots=e_{r-1}=1$, i.e.\ iff the gcd of the $(r-1)\times(r-1)$ minors of $P$ is $1$. (Known.)
\end{dimcard}

\begin{proof}
\step{1}{SYS-R1}
Every component of $u$ lies in a finite-dimensional derivative-invariant space of $\EP$ (A2), so $M_u\cong\sum Ru_i\subseteq\mathcal A$ has finite dimension over $\C$.

\step{2}{SYS-T0}
By the Smith form $M_P\cong\bigoplus_{i\le r}\C[s]/(e_i)\oplus\C[s]^{k-r}$, and $\dim\mathrm{soc}_{(s-\lambda)}(M_P)=\#\{i:e_i(\lambda)=0\}$ (a cyclic torsion summand with root $\lambda$ has a one-dimensional socle there; free summands have none). Since $e_i\mid e_r$, ``$\le1$ for every $\lambda$'' means $e_{r-1}$ has no root, i.e.\ $e_1=\dots=e_{r-1}=1$; the gcd of the $(r-1)$-minors is $e_1\cdots e_{r-1}$. Apply A32.

\whyworks{in one variable the socle count is read off the Smith form.}
\end{proof}

\begin{dimcard}{B74}{SYS-E}{Dirac $1+1$ as an evolution: what survives of the no-product theorem}{\proved}{prop:syse}
Write $\psi_t=A(\partial_x)\psi$, $A(\xi)=X+mY$, $X=i\xi\sigma_z$ (transport), $Y=-i\sigma_x$ (mass), in the algebra $\mathrm{Mat}_2(\C[\xi])$.
\begin{enumerate}[label=(E\arabic*)]
\item $X,Y_0=Y,Y_k=\mathrm{ad}_X^kY$ are linearly independent over $\C$. Hence (P1-T5, A24) no finite word $(Z_1,\dots,Z_d)\in\{X,Y\}^d$ and holomorphic $g_i(t,m)$ with $g_i(0,m)=0$, $g_i(t,0)\equiv0$ for $Z_i=Y$, gives $e^{t(X+mY)}=\prod e^{g_iZ_i}$ in $\mathrm{Mat}_2(\C[\xi])[[t]]$ for $m$ near $0$.
\item $\mathrm{Lie}_\C\langle X,Y\rangle=\mathfrak L^\theta=\{Z\in\mathfrak{sl}_2\otimes\C[\xi]:\theta(Z)=Z\}$, $\theta(Z)(\xi)=\sigma_xZ(-\xi)\sigma_x$, of infinite dimension.
\item With $\omega=\sqrt{\xi^2+m^2}$, $b=-\arcsin(m\sin(\omega t)/\omega)$ and $a=\frac12\arg(\cos\omega t+i\xi\sin(\omega t)/\omega)$ one has $e^{tA}=e^{g_1X}e^{g_2mY}e^{g_1X}$ with $g_1=a/\xi$, $g_2=-b/m$: explicit, analytic in $\xi$ for $|t|<\pi/(2|m|)$, but $g_2$ is not a polynomial in $\xi$ (a bounded polynomial would be constant, while $b\to0$ at infinity and $b(t,0)=-mt\ne0$), and $e^{g_2(t,\partial_x)mY}$ is a non-local Fourier multiplier. At $t=\pi/(2|m|)$, $g_1\to\pm\infty$ as $\xi\to0^\pm$.
\end{enumerate}
\end{dimcard}

\begin{proof}
\step{1}{(E1) Independence}
With $[\sigma_z,\sigma_x]=2i\sigma_y$, $[\sigma_z,\sigma_y]=-2i\sigma_x$: $\mathrm{ad}_X\sigma_x=-2\xi\sigma_y$, $\mathrm{ad}_X\sigma_y=2\xi\sigma_x$, so $Y_k=c_k\xi^k\sigma_x$ ($k$ even) or $c_k\xi^k\sigma_y$ ($k$ odd), $c_k\ne0$; distinct degrees in $\xi$ make them independent, and $X$ has a $\sigma_z$ component absent from all $Y_k$.

\step{2}{(E2) The algebra}
$\theta$ is a Lie algebra automorphism, so $\mathfrak L^\theta$ is a subalgebra containing $X,Y$. Conversely $\mathrm{Lie}\langle X,Y\rangle$ contains $\xi^{2k}\sigma_x$, $\xi^{2k+1}\sigma_y$ (E1) and $[\sigma_x,\xi^{2k+1}\sigma_y]=2i\xi^{2k+1}\sigma_z$, a basis of $\mathfrak L^\theta$. So P1-T4 (global Wei--Norman with $\dim\mathfrak g$ functions) does not apply over $\C$.

\step{3}{(E3) Euler angles}
$A^2=-\omega^2I$ gives $e^{tA}=\cos\omega t\,I+\frac{\sin\omega t}\omega A$. The product $e^{ia\sigma_z}e^{ib\sigma_x}e^{ia\sigma_z}$ has entries $(11)=\cos b\,e^{2ia}$, $(12)=(21)=i\sin b$, $(22)=\cos b\,e^{-2ia}$; matching $(12)=-im\sin(\omega t)/\omega$ gives $\sin b=-m\sin(\omega t)/\omega$ and $|\cos\omega t+i\xi\sin(\omega t)/\omega|^2=\cos^2b$ gives $a$. The multiplier $\cos b\,I+i\sin b\,\sigma_x$ is smooth, bounded, tends to $I$ at infinity and is not constant; a convolution operator is local iff its symbol is a polynomial (Schwartz).

\whyworks{the Lie algebra generated by transport and mass is infinite-dimensional, so a uniform finite product with polynomial times is impossible; the explicit product that exists uses non-local multipliers.}
\end{proof}
\uses{A24, A32.}

\begin{caution}[title=Open (resolved in Phase 10 for real times)]
Whether finite products with times polynomial in $\xi$ exist (Q-S8SYS-1) was closed in Phase~10: no finite word at polynomial real times (SYS-DIRAC-POL); complex times remain open (Chapter~\ref{ch:cIV4}).
\end{caution}

\begin{dimcard}{B75}{LAME-3}{Lam\'e in space: an identifying displacement field and its test}{\proved}{prop:lame3}
Let $P=\rho\tau^2I_3-\mu|\xi|^2I_3-(\lambda+\mu)\xi\xi^{\top}$ with $\rho,\mu\ne0$, $\lambda+\mu\ne0$, $\lambda+2\mu\ne0$, $q_P=\rho\tau^2-(\lambda+2\mu)|\xi|^2$, $q_S=\rho\tau^2-\mu|\xi|^2$. Then $\det P=q_Pq_S^2$, $\mathrm{Ass}(M_P)=\{(q_P),(q_S)\}$ with Bass multiplicities $(1,2)$; an identifying field always exists (A33); and $u$ identifies Lam\'e iff (a) $\Ann(\mathrm{div}\,u)=(q_P)$ and (b) every relation $\sum r_i(\partial)(\mathrm{rot}\,u)_i=0$ has $r\in q_SR^3+R\xi$ up to a factor outside $(q_S)$.
\end{dimcard}

\begin{proof}
\step{1}{Structure}
$\det P=q_Pq_S^2$; $q_P,q_S$ are irreducible quadrics of rank $4$; $\bar P$ has rank $1$ mod $q_S$ and rank $2$ mod $q_P$, so over the DVR $R_{(q_S)}$ the Smith form gives $M_{(q_S)}\cong\kappa(q_S)^2$ and $M_{(q_P)}\cong\kappa(q_P)$.

\step{2}{An essential submodule}
Let $E$ be generated by the row $d=\xi^{\top}$ (div) and the rows $c_i$ of $[\xi]_\times$ (rot). Then $\xi^{\top}P=q_P\xi^{\top}$, $[\xi]_\times P=q_S[\xi]_\times$, and Helmholtz $|\xi|^2I=\xi\xi^{\top}-[\xi]_\times^2$ gives $|\xi|^2M_P\subseteq E$; $|\xi|^2\notin(q_P)\cup(q_S)$ is a non-zero-divisor, so $E$ is essential. $Rd\cap\sum Rc_i=0$ (a common element is killed by $(q_P,q_S)$ of height $2$).

\step{3}{The test}
By SYS-L0 on $E$, $u$ identifies iff (a) and (b) hold. For (b), $M_{(q_S)}=\kappa^{1\times3}/\kappa\xi^{\top}$ and $\bar r^{\top}[\xi]_\times\in\kappa\xi\iff\bar r\times\xi=0\iff\bar r\in\kappa\xi$.

\whyworks{one field identifies Lam\'e iff its P-part (div) and its S-part (the two transverse polarisations of rot) satisfy no extra equation. A pure P-wave ($\mathrm{rot}\,u=0$) never identifies.}
\end{proof}
\uses{A33, A32.}

\begin{dimcard}{B76}{G-T2}{Gauge covariance, and what gauge-invariant data can reveal}{\proved}{prop:gauge}
\textbf{(G-L0)} $e^{-i\chi}(\partial_\mu-i(A_\mu+\partial_\mu\chi))e^{i\chi}=\partial_\mu-iA_\mu$; so $L_{A+d\chi}(e^{i\chi}u)=e^{i\chi}L_Au$ for every operator polynomial in $\nabla_\mu=\partial_\mu-iA_\mu$ with constant matrix coefficients. For $\chi=c\cdot x$, $\Ann(e^{ic\cdot x}u)=\sigma_c\Ann(u)$ with $\sigma_c(r)(\zeta)=r(\zeta-ic)$.

\textbf{(G-T2)} Let $\Omega\subseteq\R^n$ be open connected and gauge group the smooth maps $g:\Omega\to U(1)$ acting as $A\mapsto A+\frac1ig^{-1}dg$. (a) $A'\sim A\iff d(A'-A)=0$ and $\oint_\gamma(A'-A)\in2\pi\Z$ for every cycle $\gamma$; on simply connected $\Omega$, $A'\sim A\iff dA'=dA$. (b) If an observation $O$ satisfies $O(gu)=O(u)$, any procedure reconstructing $A$ from $O(u)$ recovers at most the class $[A]$.
\end{dimcard}

\begin{proof}
\step{1}{(G-L0)}
$(\partial-iA-i\partial\chi)(e^{i\chi}f)=e^{i\chi}(\partial f-iAf)$.

\step{2}{(G-T2)(a)}
($\Rightarrow$) Locally $g=e^{i\chi}$ and $A'-A=d\chi$ is closed; along $\gamma$, $\frac1i\oint g^{-1}dg=2\pi\deg(g\circ\gamma)\in2\pi\Z$. ($\Leftarrow$) $B=A'-A$ is closed with periods in $2\pi\Z$, so $g(x)=\exp(i\int_{x_0}^xB)$ is path-independent, smooth, and $\frac1ig^{-1}dg=B$.

\step{3}{(G-T2)(b)}
If $L_Au=0$ and $A'=A+\frac1ig^{-1}dg$ then $L_{A'}(gu)=gL_Au=0$ and $O(gu)=O(u)$: the data do not distinguish $A$ from $A'$. On a non-simply connected $\Omega$ the holonomies are the extra datum (Aharonov--Bohm; \cf).

\whyworks{gauge acts on $A$ and on $u$ compatibly; an observation insensitive to the phase cannot see the compensating change in $A$.}
\end{proof}
\uses{none.}

\begin{dimcard}{B77}{SYS-N3}{Magnetic Schr\"odinger: one field does not identify $(A,V)$, even modulo gauge}{\proved}{prop:sysn3}
Let $n\ge2$, $u\in C^\infty(J\times\Omega)$ nowhere zero with $iu_t=-(\nabla-iA)^2u+Vu$, $A,V$ real smooth; $\rho=|u|$, $\nabla S=\mathrm{Im}(\nabla u/u)$, $S_t=\mathrm{Im}(u_t/u)$, $w=\nabla S-A$, $j=\rho^2w$. Then:
\begin{enumerate}[label=(\alph*)]
\item a real smooth pair $(A',V')$ satisfies the same equation with the same $u$ iff $A'=\nabla S-j'/\rho^2$, $V'=-S_t+\Delta\rho/\rho-|j'|^2/\rho^4$, $\mathrm{div}_xj'=-\tfrac12\partial_t(\rho^2)$;
\item the compatible pairs form the affine space $(A,V)+\{(-k/\rho^2,(|j|^2-|j+k|^2)/\rho^4):\mathrm{div}_xk=0\}$, of infinite dimension;
\item there are $k$ with $d(k/\rho^2)\ne0$: the magnetic fields differ, the pairs are not gauge-equivalent, and a gauge fixing $u$ is trivial ($gu=u\Rightarrow g=1$).
\end{enumerate}
The variations $B=A'-A=-k/\rho^2$ satisfy $\mathrm{div}(\rho^2B)=0$ (not $\mathrm{div}B=0$).
\end{dimcard}

\begin{proof}
\step{1}{(a) Madelung form}
With $u=\rho e^{iS}$, $(\nabla-iA)^2u=e^{iS}[\Delta\rho-\rho|w|^2+\frac i\rho\mathrm{div}(\rho^2w)]$ and $iu_t=e^{iS}(i\rho_t-\rho S_t)$. Real and imaginary parts give $V=-S_t+\Delta\rho/\rho-|w|^2$ and $\partial_t\rho^2+2\mathrm{div}(\rho^2w)=0$; $A'$ enters only through $w'=\nabla S-A'$.

\step{2}{(b), (c)}
Two admissible $j,j'$ differ by $k$ with $\mathrm{div}_xk=0$, and every such $k$ gives an admissible pair. For $n=2$ take $k=\nabla^\perp\psi$, $\psi\in C_c^\infty\setminus0$; $\mathrm{rot}(k/\rho^2)=\mathrm{div}(\rho^{-2}\nabla\psi)$ is not identically zero, otherwise $\psi$ would be a compactly supported solution of an elliptic divergence-form equation, so $\psi\equiv0$ by the maximum principle. For $n\ge3$ take $k=\mathrm{rot}(\psi e_3)$.

\whyworks{one complex scalar equation cannot determine $n+1$ real unknowns: the count, not gauge, is the obstruction.}
\end{proof}
\uses{B76.}

\begin{dimcard}{B78}{VAR-P5}{Dirac $1+1$ with a potential: finite speed leaves places the data never visited}{\proved}{prop:varp5}
Let $\Omega=I\times J$, $\psi\in C^1(\Omega;\C^2)$, $i\psi_t=-i\sigma_1\psi_x+m\sigma_3\psi+V(x)\psi$, residual $r=i\psi_t+i\sigma_1\psi_x-m\sigma_3\psi=V\psi$, $Y_0=\{x:\psi(t,x)=0\ \forall t\}$.
\begin{enumerate}[label=(\alph*)]
\item (scalar continuous $V$) $V=\langle\psi,r\rangle/|\psi|^2$ where $\psi\ne0$, with compatibility $r\parallel\psi$; $V$ is identified iff $\mathrm{int}\,Y_0=\emptyset$.
\item Unlike Schr\"odinger, $\mathrm{int}\,Y_0\ne\emptyset$ with $\psi\not\equiv0$ is possible: for $m=V=0$, $\Omega=(0,1)\times(0,2)$, $\psi=f(x-t)v_+$ with $\sigma_1v_+=v_+$ and $f\in C_c^\infty((0.5,0.7))\ne0$, $Y_0\supset(0,0.5]\cup[1.7,2)$ and every $V$ supported in $(0,0.5)$ is compatible.
\item (matrix $V$, pointwise class) $V(x_0)$ is determined iff $\mathrm{span}\{\psi(t,x_0):t\in I\}=\C^2$; if the span is a line $\C v$ the ambiguity is $\{W:Wv=0\}$.
\end{enumerate}
\end{dimcard}

\begin{proof}
\step{1}{(a)}
As in A36(a): $(V-W)\psi=0\iff V=W$ off $Y_0$; for continuous $V$, iff on the dense set $J\setminus Y_0$, iff $\mathrm{int}\,Y_0=\emptyset$.

\step{2}{(b)}
With $m=V=0$ the equation is $\psi_t=-\sigma_1\psi_x$, i.e.\ $\psi_{\pm,t}=\mp\psi_{\pm,x}$ on the components $\sigma_1=\pm1$, solved by $f(x-t)$. Then $\psi\ne0$ only if $x-t\in(0.5,0.7)$, which never occurs for $x\le0.5$ or $x\ge1.7$, $t\in(0,1)$.

\step{3}{(c)}
$W\mapsto(W\psi(t,x_0))_t$ is injective iff no $W\ne0$ vanishes on the span. A Hermitian $W$ with $Wv=0$ preserves $v^\perp$, so $W=cP_{v^\perp}$ with $c$ real.

\whyworks{finite propagation speed creates regions of space-time the solution never visits, where the potential cannot be read; Schr\"odinger's infinite speed, through A37, rules this out.}
\end{proof}
\uses{A36, A37.}

\begin{dimcard}{B79}{ENC-K0}{A constant potential has no strict differential intertwiner}{\proved}{prop:enck0}
Let $J$ be an open interval, $c\in\C\setminus0$ and $W=\sum_{j=0}^nw_j\partial^j$, $w_j\in C^\infty(J;\C)$, $w_n\not\equiv0$. Then $WL_c\ne L_0W$. (With the relaxed requirement $WL_c=(L_0+c)W$, $W=I$ works.)
\end{dimcard}

\begin{proof}
\step{1}{Read the coefficients}
$WL_c-L_0W=[\partial^2,W]+cW=\sum_j(w_j''\partial^j+2w_j'\partial^{j+1})+cW$. If it vanishes: the coefficient of $\partial^{n+1}$ gives $2w_n'=0$, so $w_n$ is a non-zero constant; for $1\le j\le n$, $w_j''+2w_{j-1}'+cw_j=0$; and $w_0''+cw_0=0$.

\step{2}{Induction on $k=0,\dots,n$}
$w_{n-k}$ is a polynomial of degree exactly $k$ with leading coefficient $(-c/2)^kw_n/k!$: if $w_j$ has degree $k$ and leading coefficient $\ell$, then $w_{j-1}'=-(w_j''+cw_j)/2$ has degree $k$ and leading coefficient $-c\ell/2$, so $w_{j-1}$ has degree $k+1$ and leading coefficient $-c\ell/(2(k+1))$.

\step{3}{Contradiction}
For $k=n$, $w_0$ has degree $n$ and non-zero leading coefficient, so $w_0''+cw_0$ has degree $n$ and leading coefficient $c(-c/2)^nw_n/n!\ne0$, against the last equation.

\whyworks{the strict equation forces a polynomial tower that cannot close; the constant $c$ must be moved to the right-hand side.}
\end{proof}
\uses{none.}

\begin{dimcard}{B149}{SYS-T1}{An identification criterion on an essential family}{\pw{} (a tool, not a structural theorem)}{ch:systems}
Let $M$ be a finitely generated $R$-module and $\{x_{\mathfrak p,j}\}$ a finite family with $Rx_{\mathfrak p,j}\cong R/\mathfrak p$, $\mathfrak p\in\mathrm{Ass}(M)$, whose sum $E$ is direct and essential in $M$. For $u\in\mathrm{Sol}_{\mathcal A}(P)$ put $w_{\mathfrak p,j}=\varphi_u(x_{\mathfrak p,j})$. Then $u$ identifies $P$ iff for every $\mathfrak q\in\mathrm{Ass}(M_P)$ there is no relation $\sum_{\mathfrak p\supseteq\mathfrak q}\sum_jr_{\mathfrak p,j}(\partial)w_{\mathfrak p,j}=0$ with some $r_{\mathfrak q,j}\notin\mathfrak q$.
\end{dimcard}

\begin{proof}
\step{1}{Injectivity can be tested on an essential submodule (SYS-L0)}
Let $\varphi:M\to F$ be a homomorphism and $E\subseteq M$ essential. Then $\ker\varphi\cap E=\ker(\varphi|_E)$; if this is $0$, then $\ker\varphi$ meets $E$ trivially, hence is zero by essentiality. So $\varphi$ is injective iff $\varphi|_E$ is.

\step{2}{A non-zero kernel element has prime annihilator}
A non-zero element of $\ker\varphi_u\cap E$ can be chosen with annihilator a prime $\mathfrak q\in\mathrm{Ass}(E)=\mathrm{Ass}(M_P)$. Its expression involves only summands with $\mathfrak p\supseteq\mathfrak q$ and at least one coefficient $r_{\mathfrak q,j}\notin\mathfrak q$: a finite intersection of primes equal to a prime $\mathfrak q$ contains the product of those primes, hence $\mathfrak q$ contains one of them.

\step{3}{Reading}
So $\varphi_u$ fails to be injective iff such a relation exists. For $k=1$ this is exactly C3-T1.

\whyworks{a module map that kills nothing on an essential submodule kills nothing at all.}
\end{proof}
\uses{Matlis (essential extensions), cited.}

\begin{caution}[title={Status and prior art}]
SYS-T1 is not a structural theorem: it restates SYS-L0 and is no more decidable than computing $\mathrm{Ann}(u)$ directly for non-EP $u$. It shortens the test to $\mu_{\mathfrak p}$ scalar functions per associated prime (Dirac $3{+}1$, Lam\'e). The structural content is SYS-EX. For exponential-polynomial data $M_u$ has finite length (SYS-R1), the setting of Schindelar--Levandovskyy--Zerz 2011 (abstract read only).
\end{caution}

\begin{dimcard}{B150}{SYS-N1}{The socle obstruction}{\pw}{ch:systems}
Let $\Omega$ be connected and $M$ an $R$-module. If $\dim_{\C}\operatorname{soc}_{\mathfrak m_\lambda}(M)\ge2$ for some maximal ideal $\mathfrak m_\lambda$, no homomorphism $M\to\mathcal A$ is injective. In particular if this holds for $M_P$, no $u\in\mathrm{Sol}_{\mathcal A}(P)$ identifies $P$.
\end{dimcard}

\begin{proof}
\step{1}{The image of the socle}
A homomorphism $\varphi$ sends $\operatorname{soc}_{\mathfrak m_\lambda}(M)$ into $\{f\in\mathcal A:(\partial_j-\lambda_j)f=0\ \forall j\}$.

\step{2}{That space is a line}
On a connected $\Omega$, write $f=e^{\lambda\cdot x}g$: then $\nabla g=0$ in $\mathcal D'(\Omega)$, so $g$ is constant, and the space is $\C e^{\lambda\cdot x}$.

\step{3}{No injectivity}
A linear map from a space of dimension $\ge2$ to a line is not injective.

\whyworks{a connected domain carries only one exponential per eigenvalue, so a socle with two copies cannot embed.}
\end{proof}
\uses{none.}

\begin{caution}[title={Prior art}]
On $\Omega=\R^n$ SYS-N1, SYS-EX and SYS-T0 are known (Oberst 1993 for the minimal number of trajectories, Pommaret, Macaulay for the length-finite case); the only new point is arbitrary connected $\Omega$ (A32), where $C^\infty(\Omega)$ is not injective. The texts of Oberst 1990/1993 were not read.
\end{caution}

\begin{dimcard}{B151}{VAR-T6}{Order of the minimal operator, socle and cyclicity (Picard--Vessiot)}{\cf}{ch:varcoef}
Let $V$ be the solution space of $L$ in a Picard--Vessiot extension $E\supset K$ with group $G$, $u\in V\setminus0$, $W_u=\operatorname{span}_{\C}(G\cdot u)$. (a) $G$-stable subspaces $W\subset V$ correspond to monic right factors $L_W=\mathrm{Wr}(y_1..y_k,Y)/\mathrm{Wr}(y_1..y_k)$ of $L$, with $\mathrm{Sol}_E(L_W)=W$. (b) $\mathrm{Sol}_E(L_u)=W_u$, so $\operatorname{ord}L_u=\dim\operatorname{span}(G\cdot u)$ and $u$ identifies $L$ iff $u$ is a cyclic vector of $V$ over $\C[G]$. (c) An identifying $u$ exists iff $V$ is cyclic over $\mathcal A=\operatorname{span}(G)$; if $V\cong\bigoplus S_i^{m_i}$ is semisimple, iff $m_i\le\dim S_i$. (d) The right factor $C_1$ attached to $\operatorname{soc}(V)$ induces $V/\operatorname{soc}V\cong\mathrm{Sol}(L_1)$ for $L=L_1C_1$; iterating gives the Loewy series.
\end{dimcard}

\begin{proof}
\step{1}{(a)}
For $W$ $G$-stable with basis $y_i$, $g(y_1..y_k)=(y_1..y_k)A_g$, both Wronskians are multiplied by $\det A_g$, so the coefficients of $L_W$ are $G$-invariant, i.e.\ in $K$ ($E^G=K$, cited). $\mathrm{Sol}(L_W)\supset W$ and has dimension $\le k$. Division $L=QL_W+R$ with $\operatorname{ord}R<k$ gives $R|_W=0$, so $R=0$. Conversely, if $L=QR$, $\mathrm{Sol}_E(R)\subset V$ is $G$-stable of dimension $\operatorname{ord}R$ (since $R$ maps $V$ into $\mathrm{Sol}(Q)$).

\step{2}{(b)}
$\mathrm{Sol}(L_u)$ is $G$-stable and contains $u$, hence $\supset W_u$; conversely $L_{W_u}$ kills $u$, so $L_{W_u}\in\mathcal DL_u$ and $\mathrm{Sol}(L_u)\subset W_u$. Thus $u$ identifies iff $W_u=V$.

\step{3}{(c)}
The first statement is (b). If $V$ is semisimple, $\mathcal A$ acts through $\bigoplus\mathrm{End}(S_i)$ (Jacobson density); $S_i^{m}\cong S_i\otimes\C^m$ is cyclic over $\mathrm{End}(S_i)$ iff $m\le\dim S_i$ (a vector is cyclic iff its tensor rank is $m$), and components at distinct factors are independent.

\step{4}{(d)}
$C_1$ commutes with $G$ (coefficients in $K$), so $C_1:V\to E$ is a $G$-morphism with kernel $\mathrm{Sol}(C_1)=\operatorname{soc}V$ by (a); its image lies in $\mathrm{Sol}(L_1)$ and has dimension $n-\operatorname{ord}C_1=\operatorname{ord}L_1$.

\whyworks{a solution identifies the operator exactly when its Galois orbit spans the whole solution space.}
\end{proof}
\uses{van der Put--Singer ($E^G=K$, numbering to check).}

\begin{caution}[title={Examples that follow}]
Airy $\partial^2-x$ has no rational Riccati solution, so every non-zero solution identifies it. For $L_E=-\partial^2+x^2-E$ the rational solutions are $p(x)e^{\mp x^2/2}$ with $p$ Hermite, existing iff $E=\pm(2n+1)$; for $E\notin\pm(2\N+1)$ every $u\ne0$ identifies $L_E$, and for $E=2n+1$ one has $u$ identifying iff $u\notin\C h_n$ (Loewy length~2, socle $\C h_n$).
\end{caution}

\section{Statements recorded without proof here}

\listcard{SYS-T1}{Identification criterion on an essential family: proved in B149 (as a tool).}{\pw{} (a tool, not a structural theorem)}{S8-SYS \S1. Reduces to C3-T1 for $k=1$; the structural content is SYS-EX.}
\listcard{SYS-N1, SYS-EX, SYS-T0}{Socle obstruction (proved in B150); existence criterion; ODE systems on $\Omega=\R^n$.}{known (Oberst, Pommaret, Macaulay)}{The only new point is arbitrary $\Omega$ (A32).}
\listcard{SYS-S1-C}{Failure of (i$\Leftarrow$) on a simply connected non-convex $\Omega$, in $C^\infty$ and $\mathcal D'$.}{\falsest}{Proved in Chapter~\ref{ch:cIII1}, after A34.}
\listcard{SYS-II-T2-C}{``Crossing of eigenvalues implies a pole of the projector'' is false for a derogatory $A$ ($A=\mathrm{diag}(\xi,-\xi)$ has constant projectors); see B138.}{\falsest}{one-line counterexample, S8-SYS \S8.}
\listcard{Q-S8SYS-4}{``The pole depends only on $A(\xi_0)$ and the factorisation''; the counterexample pair $A_1,A_2$ is in B138.}{\falsest}{Phase~10 (F10-C).}
\listcard{G-T0}{A gauge-fixed Maxwell potential never identifies Maxwell by itself; if $\Ann(A)=U_A+Rg$ with $g\cdot\partial\ne0$ then $U_A=\{r\in\Ann(A):r\cdot\partial=0\}$.}{observation}{S8-SYS \S6.}

\section{Variable coefficients}

\listcard{VAR-T6}{Order of $L_u$ is $\dim\mathrm{span}(Gu)$; socle/Loewy structure; existence iff $V$ is cyclic ($m_i\le\dim S_i$). Proved from $E^G=K$ in B151.}{from cited source}{Galois correspondence $E^G=K$; S8-VAR \S2.2.}
\listcard{VAR-P3 (Lipschitz, no continuity from $C^1$, H\"older $1/3$ sharp)}{See A39.}{\pw}{Chapter~\ref{ch:cIII2}.}

\noindent{\small \textbf{Proofs moved.} The results VAR-T2a, VAR-E1/E2, VAR-P8, VAR-P9, VAR-P11, VAR-T10, VAR-P12 and VAR-C3Q2 have a proof card in Chapter~\ref{ch:cIV5} (Dim.~B80--B86). VAR-T6 (cited) and VAR-P3 stay as above.}
